\subsection{补遗二：生成}

\ssszh 最后一个话题是几条关于生成的好定理。回忆一条基本而有用的常识：Morita 理论告诉我们，给定环 $A$，其模范畴 $A\mod$ 虽由 $A$ 本身紧生成，却也容许许多别的紧生成元 $c$。任取其一，取它的自同态就给出等价
%
\[
  \rmod\End(c)  \simeq   A\mod, \quad \quad M \mapsto M \otimes_{\End(c)} c.
\]
%
换一种说法，给定的 dg 范畴 $\sC$ 可以有许多紧生成元，相应地也有许多把它表示成模范畴的方式。这里想综述两条有用的此类陈述；它们必然发生在高一个范畴层级上，即对 $\Dmod(K)\mod$ 而言。

\ssszh \label{ss:hchbimod}首先，Beraldo \cite{beraldo} 在 Gaitsgory \cite{1affine} 的工作基础上证明了一条非平凡的定理：$\Dmod(K)$ 由 $\fk\mod$ 生成。注意后者的自同态为
%
\begin{align*}
   \Hom_{\Dmod(K)\mod}(\fk\mod, \fk\mod) & \simeq \udHom(\fk\mod, \fk\mod)^K \\
   & \simeq (\fk\mod^\vee \otimes \fk\mod)^K \\ & \simeq (\fk\mod \otimes
   \fk\mod)^K \\ & \simeq (\fk \oplus \fk\mod)^K.
\end{align*}
%
用文字表述：$\fk\mod$ 的等变自同态构成 Harish-Chandra 双模范畴，也就是对角作用可积为群 $K$ 的作用的 $\fk$-双模。\zy{[译者按（上下文推断，无外部依据）：按上面最后的计算，此行应读作 $(\fk \oplus \fk)\mod$，即 $\fk \oplus \fk$ 的模范畴；原文把它印作 $\fk \oplus \fk\mod$，本译本照原文排印。]}

\begin{re} 我们要强调，这条等同是该形式体系一个漂亮而有用的特点。也就是说，Harish-Chandra 双模在表示论中已被研究和使用数十年，人们早已认识到它们具有许多良好的性质，参见 \cite{bege1980} 及其中的参考文献。上面的陈述在精确的意义上说明了这一点：它们恰好就是与范畴 $K$-作用交换的 $\fk\mod$ 的自函子，反之亦然。
\end{re}

于是 Beraldo 的定理给出等价
%
\[
    \Dmod(K)\mod \simeq (\fk \oplus \fk\mod^K)\mod.
\]
%
特别地，对任意 $\Dmod(K)$ 表示 $\sC$，有典范等价
%
\[
    \fk\mod \otimes_{\fk\oplus \fk\mod^K} \sC^{K, w} \simeq \sC.
\]
%
也就是说，回忆等同 $\Dmod(K)^{K,w} \simeq \fk\mod$：一个 $\Dmod(K)$ 表示可以由它的弱不变量重构出来。

\ssszh 对约化群 $G$，我们还想提一条由 Ben-Zvi--Gunningham--Orem \cite{BZO} 给出的最高权理论在当下框架中的漂亮类比。他们证明了：对任意极大幂幺子群 $N$，任意的 $\Dmod(G)$ 表示 $\sC$ 都可以由它的 $N$-不变量重构，甚至只需后者的弱 Cartan 不变量：
%
\[
   \sC^N \quad \text{和} \quad (\sC^N)^{T,w}.
\]
%
过渡到它们的自同态（这些自同态大致是沿 $T$-轨道没有指定单值性的 Hecke 范畴），便得到等价
%
\[
  \Dmod(G)\mod \simeq \Dmod(N \backslash G / N)\mod \simeq \Dmod(N \backslash G
  / N)^{T \times T, w}\mod
\]
%
特别地，还得到等同
%
\[
   \Dmod(G/N) \otimes_{\Dmod(N \backslash G / N)}  \sC^N \simeq
   \Dmod(G/N)^{T,w} \otimes_{\Dmod(N \backslash G / N)^{T \times T,w}}
   (\sC^{N})^{T,w} \simeq \sC.
\]
%


\begin{re} \cite{BZO} 的论证还给出下面这条有用的变体。给定群 $K$ 的子群 $H$，使得 $K$ 对 $H$ 的正规化子的商是真紧的，则典范映射
	%
	\[
	     \Dmod(K/H) \otimes_{\Dmod(H \backslash K / H)} \sC^K
	     \rightarrow \sC,
	\]
	%
虽未必本质满，却总是全忠实的嵌入。由此可以检验「最高权向量」上的范畴等价，而后者往往容易直接分析得多；进一步的讨论与例子参见 \cite{ahc}、\cite{yangmoyprasad}。
\end{re}
